%======================================================================
%  CRESCO8 as the experimental testbed
%======================================================================

\section{CRESCO8 as the experimental testbed}
\label{sec:cresco8}

\subsection{Documented configuration}

CRESCO8 is the current CRESCO system at the ENEA research centre in
Portici.  It entered service in 2025 and is connected by a 200\,Gb/s
Mellanox NDR fabric~\cite{cresco8doc,eneapress}.  The operational user
page distinguishes three compute-node classes:

\begin{table}[tbp]
\centering
\small
\caption{CRESCO8 node classes documented by ENEAGRID.  Counts describe a
dated operational document, not a permanent scheduler inventory
~\cite{cresco8doc}.}
\label{tab:cresco8-node-classes}
\begin{tabularx}{\textwidth}{@{}lcccY@{}}
\toprule
\textbf{Class} & \textbf{Nodes} & \textbf{CPU cores/node}
 & \textbf{Memory/node} & \textbf{Processor or accelerator} \\
\midrule
CPU
 & 760 & 128 & 512\,GB DDR5
 & 2$\times$ Xeon Platinum 8592+ \\
CPU--HBM
 & 15 & 112 & 1,024\,GB DDR5 + 128\,GB HBM
 & 2$\times$ Xeon CPU Max 9480 \\
GPU
 & 17 & 128 & 512\,GB DDR5 + accelerator HBM
 & 2$\times$ Xeon Platinum 8592+ and
   4$\times$ Intel Data Center GPU Max 1550 \\
\bottomrule
\end{tabularx}
\end{table}

The same page calls this a 793-node system, although the listed classes
sum to 792 and its reported core total is consistent with the 792-node
breakdown.  Earlier commissioning material gives yet another snapshot
~\cite{eneacresco8slides}.  This inconsistency is operational rather
than scientific: every experiment should retain dated
\texttt{sinfo -N} and \texttt{scontrol show node} output, and the thesis
should report the observed inventory rather than silently reconcile
historical documents.

\FloatBarrier

\subsection{Why the HBM partition supports a clean comparison}

Each documented HBM node contains two Max 9480 sockets, 1\,TB of DDR5
and 128\,GB of HBM.  This makes three comparisons possible, but they do
not have equal evidential strength:

\begin{enumerate}
  \item \textbf{Local HBM versus local DDR5 on the same Max 9480
        socket} is the primary comparison.  Processor, core count,
        executable and most firmware conditions can be held constant.

  \item \textbf{One socket versus two sockets, or Quadrant versus SNC4},
        tests the runtime's response to topology.  It is a placement
        experiment, not a pure memory-technology comparison.

  \item \textbf{Max 9480 versus Platinum 8592+} is a secondary
        system-level baseline.  The processors differ in generation,
        core count, cache, clocks, memory channels and uncore behaviour.
        The contrast cannot identify an HBM effect by itself.
\end{enumerate}

The installed DDR capacity does not reveal the DIMM layout.  Sixteen
64\,GB DIMMs would populate all channels, but the platform also permits
eight 128\,GB DIMMs with Xeon Max~\cite{lp1603}.  Channel population,
memory frequency and local bandwidth therefore require direct
inspection.  The presence of DDR5 also means that the nodes use Flat or
Cache mode rather than the no-DDR HBM-only configuration, but the ENEA
page does not identify which mode is active.  In Flat mode, HBM appears
as memory-only NUMA nodes; in Cache mode it does not appear separately.

The NDR link has a raw one-direction rate of 25\,GB/s per node, far
below local HBM bandwidth.  Multi-node inference may still be necessary
for capacity, but communication can then dominate the placement effect.
Single-node experiments should establish a stable result before adding
distributed parallelism.

\subsection{Scheduler and storage context}

The documented production partition is \texttt{cresco8\_hbm}, with 14
nodes and a 24-hour limit; \texttt{cresco8\_hbm\_dbg} lists two nodes
and a one-hour limit~\cite{cresco8doc}.  Their hostnames can overlap, so
these counts must not be added.

Compute nodes are reached only through Slurm.  The front ends are for
editing, compilation and submission, not benchmarking.  ENEAGRID
provides AFS for shared home directories and Lustre for parallel I/O;
large model files and benchmark output belong on the latter, subject to
the site's current mount points and quotas.  Compiler, MPI and profiler
modules listed on the web page are useful starting points, but an
experiment must record the modules and linked libraries actually loaded
on the run date.

\subsection{First allocation: capture facts before benchmarking}
\label{sec:first-allocation}

The first HBM allocation should create a machine-readable system
manifest.  The following script is deliberately a topology-capture
skeleton rather than a performance benchmark; account names, available
utilities and Slurm resource syntax must be checked on the live system.

\begin{lstlisting}[style=shell,
  caption={Minimal exclusive allocation for a CRESCO8 topology
  manifest.}]
#!/bin/bash
#SBATCH --nodes=1
#SBATCH --ntasks=1
#SBATCH --cpus-per-task=112
#SBATCH --partition=cresco8_hbm
#SBATCH --exclusive
#SBATCH --hint=nomultithread
#SBATCH --output=%j-system.txt

set -u

echo "job_id=${SLURM_JOB_ID}"
echo "host=$(hostname)"
date --iso-8601=seconds
scontrol show job "${SLURM_JOB_ID}"
lscpu
lscpu --extended
numactl --hardware
lstopo-no-graphics 2>/dev/null || true
cat /sys/devices/system/cpu/smt/active 2>/dev/null || true
findmnt -t lustre 2>/dev/null || true
module list 2>&1 || true
\end{lstlisting}

The manifest should be extended, subject to permissions, with the
following evidence:

\begin{table}[tbp]
\centering
\small
\caption{Minimum platform record for an inference result.}
\label{tab:platform-record}
\begin{tabularx}{\textwidth}{@{}lY@{}}
\toprule
\textbf{Category} & \textbf{Record} \\
\midrule
Identity
 & Hostname, Slurm job ID, partition, date and node allocation. \\
Processor
 & Exact SKU, sockets, physical cores, SMT state, microcode, governor,
   turbo policy and observed frequency. \\
Memory
 & Memory and clustering modes, NUMA distance matrix, HBM and DDR
   node IDs, DIMM population, huge-page state and measured local/remote
   bandwidth. \\
I/O
 & Lustre mount and striping used for model loading; NIC identity and
   NUMA affinity if multi-node tests are attempted. \\
Software
 & OS and kernel, loaded modules, compiler, runtime commit, linked
   libraries, build flags and complete command line. \\
Telemetry
 & Available RAPL, performance-counter and management-controller
   interfaces, including permissions and measurement boundaries. \\
\bottomrule
\end{tabularx}
\end{table}

This separation between documented configuration and observed
configuration is important.  Vendor documents establish what the
platform can support, and ENEAGRID describes a dated cluster state.
Only the archived manifest establishes the conditions under which a
reported result was obtained.
